% ==============================================================================
% config/estilo.tex
% Aparência do documento: títulos, legendas, ilustrações, códigos, referências
% cruzadas, siglas e caixas de orientação. Em geral, não é preciso alterar.
% ==============================================================================

% ------------------------------------------------------------------------------
% Termo usado para as divisões principais do texto
% ------------------------------------------------------------------------------
% A NBR 14724:2024 (item 4.2.2) determina que o trabalho seja organizado em
% seções, e não em capítulos. O padrão abaixo segue as regras de TCC do curso.
% Para adotar a terminologia da norma, troque por {Seção} e {Seções} e revise
% os artigos no texto ("o Capítulo 2" passa a ser "a Seção 2").
\newcommand{\termocap}{Capítulo}
\newcommand{\termocaps}{Capítulos}

% ------------------------------------------------------------------------------
% Parágrafos e títulos
% ------------------------------------------------------------------------------
\setlength{\parindent}{1.25cm}

% Destaque gráfico progressivo das seções (NBR 6024). O caixa alta das seções
% primárias vem da opção chapter=TITLE da classe.
\renewcommand{\ABNTEXchapterfont}{\normalfont\bfseries}
\renewcommand{\ABNTEXchapterfontsize}{\large}
\renewcommand{\ABNTEXsectionfont}{\normalfont\bfseries}
\renewcommand{\ABNTEXsectionfontsize}{\normalsize}
\renewcommand{\ABNTEXsubsectionfont}{\normalfont\bfseries\itshape}
\renewcommand{\ABNTEXsubsectionfontsize}{\normalsize}
\renewcommand{\ABNTEXsubsubsectionfont}{\normalfont\itshape}
\renewcommand{\ABNTEXsubsubsectionfontsize}{\normalsize}

% ------------------------------------------------------------------------------
% Legendas e fontes das ilustrações e tabelas
% ------------------------------------------------------------------------------
% A NBR 14724 pede fonte menor e uniforme para legendas, fontes, notas e
% citações longas. O abnTeX2 usa \footnotesize para \fonte e citações longas.
\captionsetup{font=footnotesize}
\captionsetup[sub]{font=footnotesize}

% Texto padrão para material elaborado pelo próprio autor (NBR 14724:2024).
\newcommand{\fonteautor}{\fonte{Elaborado pelo próprio autor.}}

% Posição padrão dos objetos flutuantes.
\setfloatlocations{figure}{htbp}
\setfloatlocations{table}{htbp}

% ------------------------------------------------------------------------------
% Novos tipos de ilustração: Quadro, Código e Algoritmo
% ------------------------------------------------------------------------------
% Quadro: informação textual organizada em linhas e colunas, com bordas
% fechadas. Tabela: dados numéricos, sem bordas laterais (normas do IBGE).
\newcommand{\quadroname}{Quadro}
\newcommand{\listofquadrosname}{Lista de quadros}
\newfloat{quadro}{loq}{\quadroname}
\newlistof{listofquadros}{loq}{\listofquadrosname}
\newlistentry{quadro}{loq}{0}
\setfloatlocations{quadro}{htbp}
\renewcommand{\cftquadroname}{\quadroname\space}
\renewcommand*{\cftquadroaftersnum}{\hfill\textendash\hfill}

\newcommand{\codigoname}{Código}
\newcommand{\listofcodigosname}{Lista de códigos}
\newfloat{codigo}{lcod}{\codigoname}
\newlistof{listofcodigos}{lcod}{\listofcodigosname}
\newlistentry{codigo}{lcod}{0}
\setfloatlocations{codigo}{htbp}
\renewcommand{\cftcodigoname}{\codigoname\space}
\renewcommand*{\cftcodigoaftersnum}{\hfill\textendash\hfill}

\newcommand{\algoritmoname}{Algoritmo}
\newcommand{\listofalgoritmosname}{Lista de algoritmos}
\newfloat{algoritmo}{lalg}{\algoritmoname}
\newlistof{listofalgoritmos}{lalg}{\listofalgoritmosname}
\newlistentry{algoritmo}{lalg}{0}
\setfloatlocations{algoritmo}{htbp}
\renewcommand{\cftalgoritmoname}{\algoritmoname\space}
\renewcommand*{\cftalgoritmoaftersnum}{\hfill\textendash\hfill}

% Largura reservada para o rótulo ("Quadro 10", "Código 10" etc.) nas listas.
\setlength{\cftquadronumwidth}{2.5em}
\setlength{\cftcodigonumwidth}{2.5em}
\setlength{\cftalgoritmonumwidth}{2.5em}

% ------------------------------------------------------------------------------
% Códigos (pacote listings)
% ------------------------------------------------------------------------------
\definecolor{codigofundo}{gray}{0.96}
\definecolor{codigocomentario}{gray}{0.40}
\definecolor{codigonumero}{gray}{0.55}

\lstset{
  basicstyle=\ttfamily\footnotesize\linespread{1}\selectfont,
  keywordstyle=\bfseries,
  commentstyle=\itshape\color{codigocomentario},
  stringstyle=\color{black},
  numbers=left,
  numberstyle=\tiny\color{codigonumero},
  numbersep=8pt,
  backgroundcolor=\color{codigofundo},
  frame=tb,
  framerule=0.4pt,
  xleftmargin=1.5em,
  breaklines=true,
  breakatwhitespace=true,
  showstringspaces=false,
  tabsize=4,
  columns=fullflexible,
  keepspaces=true,
  upquote=true,
  extendedchars=true,
  inputencoding=utf8,
  % Acentos dentro de códigos e comentários (necessário com pdfLaTeX).
  literate=
    {á}{{\'a}}1 {à}{{\`a}}1 {ã}{{\~a}}1 {â}{{\^a}}1
    {é}{{\'e}}1 {ê}{{\^e}}1 {í}{{\'i}}1
    {ó}{{\'o}}1 {õ}{{\~o}}1 {ô}{{\^o}}1 {ú}{{\'u}}1 {ü}{{\"u}}1
    {ç}{{\c{c}}}1
    {Á}{{\'A}}1 {À}{{\`A}}1 {Ã}{{\~A}}1 {Â}{{\^A}}1
    {É}{{\'E}}1 {Ê}{{\^E}}1 {Í}{{\'I}}1
    {Ó}{{\'O}}1 {Õ}{{\~O}}1 {Ô}{{\^O}}1 {Ú}{{\'U}}1
    {Ç}{{\c{C}}}1
}

% ------------------------------------------------------------------------------
% Algoritmos (pacote algpseudocode) com palavras-chave em português
% ------------------------------------------------------------------------------
\algrenewcommand\algorithmicrequire{\textbf{Entrada:}}
\algrenewcommand\algorithmicensure{\textbf{Saída:}}
\algrenewcommand\algorithmicif{\textbf{se}}
\algrenewcommand\algorithmicthen{\textbf{então}}
\algrenewcommand\algorithmicelse{\textbf{senão}}
\algrenewcommand\algorithmicend{\textbf{fim}}
\algrenewcommand\algorithmicwhile{\textbf{enquanto}}
\algrenewcommand\algorithmicdo{\textbf{faça}}
\algrenewcommand\algorithmicfor{\textbf{para}}
\algrenewcommand\algorithmicforall{\textbf{para todo}}
\algrenewcommand\algorithmicrepeat{\textbf{repita}}
\algrenewcommand\algorithmicuntil{\textbf{até que}}
\algrenewcommand\algorithmicreturn{\textbf{retorne}}
\algrenewcommand\algorithmicfunction{\textbf{função}}
\algrenewcommand\algorithmicprocedure{\textbf{procedimento}}
\algrenewcommand{\algorithmiccomment}[1]{\hfill\(\triangleright\) #1}

% ------------------------------------------------------------------------------
% Gráficos (pgfplots) com vírgula decimal e ponto como separador de milhar
% ------------------------------------------------------------------------------
\pgfplotsset{
  /pgf/number format/.cd, use comma, 1000 sep={.},
}

% ------------------------------------------------------------------------------
% Referências cruzadas (cleveref)
% ------------------------------------------------------------------------------
% Use \cref{rotulo} no meio da frase e \Cref{rotulo} no início. Nos dois casos
% o nome sai com inicial maiúscula ("Figura 3", "Tabela 2"), como pede a norma.
\crefname{chapter}{\termocap}{\termocaps}
\Crefname{chapter}{\termocap}{\termocaps}
\crefname{section}{Seção}{Seções}
\Crefname{section}{Seção}{Seções}
\crefname{subsection}{Seção}{Seções}
\Crefname{subsection}{Seção}{Seções}
\crefname{subsubsection}{Seção}{Seções}
\Crefname{subsubsection}{Seção}{Seções}
\crefname{appendix}{Apêndice}{Apêndices}
\Crefname{appendix}{Apêndice}{Apêndices}
% O abnTeX2 não aciona a detecção de apêndices do cleveref. Dentro dos
% ambientes de apêndices e anexos, os capítulos passam a ser referenciados
% como "Apêndice A" e "Anexo A".
\AtBeginEnvironment{apendicesenv}{\crefalias{chapter}{appendix}}
\AtBeginEnvironment{anexosenv}{\crefalias{chapter}{anexo}}
\crefname{figure}{Figura}{Figuras}
\Crefname{figure}{Figura}{Figuras}
\crefname{subfigure}{Figura}{Figuras}
\Crefname{subfigure}{Figura}{Figuras}
\crefname{table}{Tabela}{Tabelas}
\Crefname{table}{Tabela}{Tabelas}
\crefname{quadro}{Quadro}{Quadros}
\Crefname{quadro}{Quadro}{Quadros}
\crefname{codigo}{Código}{Códigos}
\Crefname{codigo}{Código}{Códigos}
\crefname{algoritmo}{Algoritmo}{Algoritmos}
\Crefname{algoritmo}{Algoritmo}{Algoritmos}
\crefname{equation}{Equação}{Equações}
\Crefname{equation}{Equação}{Equações}
\crefname{page}{página}{páginas}
\Crefname{page}{Página}{Páginas}
\creflabelformat{equation}{(#2#1#3)}
\crefrangelabelformat{equation}{(#3#1#4) a (#5#2#6)}
\newcommand{\crefpairconjunction}{ e }
\newcommand{\crefmiddleconjunction}{, }
\newcommand{\creflastconjunction}{ e }
\newcommand{\crefrangeconjunction}{ a }

% ------------------------------------------------------------------------------
% Siglas (pacote acro)
% ------------------------------------------------------------------------------
% A lista mostra só as siglas usadas no texto, em ordem alfabética.
\acsetup{
  list/template = longtable,
  list/heading  = none,
  list/display  = used,
  list/sort     = true,
}

% ------------------------------------------------------------------------------
% Links do PDF
% ------------------------------------------------------------------------------
\definecolor{azulescuro}{RGB}{0,51,128}
\hypersetup{
  colorlinks = true,
  linkcolor  = azulescuro,
  citecolor  = azulescuro,
  urlcolor   = azulescuro,
  filecolor  = azulescuro,
  bookmarksdepth = 4,
}
% Para imprimir, links pretos ficam melhores. Descomente a linha abaixo:
% \hypersetup{hidelinks}

% ------------------------------------------------------------------------------
% Notas de revisão (aparecem só com \mostrarnotastrue no arquivo principal)
% ------------------------------------------------------------------------------
% Uso: \notaaluno{texto} e \notaorientador{texto}.
\newcommand{\notaaluno}[1]{\todo[author=\textbf{Aluno},color=green!25,inline,caption={}]{#1}}
\newcommand{\notaorientador}[1]{\todo[author=\textbf{Orientador},color=red!20,inline,caption={}]{#1}}

% ------------------------------------------------------------------------------
% Caixas de orientação e de exemplos de LaTeX
% ------------------------------------------------------------------------------
% As caixas "Orientação" explicam o que escrever em cada parte do trabalho.
% Elas somem do PDF com \mostrarinstrucoesfalse no arquivo principal.
% IMPORTANTE: \begin{instrucao} e \end{instrucao} devem ficar sozinhos na linha.
\definecolor{orientacaoborda}{RGB}{30,80,140}
\definecolor{orientacaofundo}{RGB}{240,245,252}
\definecolor{exemploborda}{gray}{0.45}

\ifmostrarinstrucoes
  \newtcolorbox{instrucao}{
    enhanced, breakable,
    colback=orientacaofundo, colframe=orientacaoborda,
    boxrule=0.6pt, arc=2pt,
    left=6pt, right=6pt, top=4pt, bottom=4pt,
    title={Orientação},
    fonttitle=\sffamily\bfseries\footnotesize,
    fontupper=\small\linespread{1}\selectfont,
    before upper={\setlength{\parindent}{0pt}\setlength{\parskip}{0.5em}},
    before skip=0.8em, after skip=0.8em,
  }
\else
  \excludecomment{instrucao}
\fi

% Subtítulo dentro de uma caixa de orientação.
\newcommand{\orientacao}[1]{{\sffamily\bfseries #1}\par}

% Código LaTeX e o resultado logo abaixo.
\newtcblisting{exemplo}{
  enhanced, listing and text,
  colback=white, colframe=exemploborda, boxrule=0.4pt, arc=1pt,
  left=4pt, right=4pt, top=2pt, bottom=2pt,
  segmentation style={dashed, exemploborda},
  title={Código LaTeX e resultado}, coltitle=black, colbacktitle=black!8,
  fonttitle=\sffamily\footnotesize,
  listing options={language={[LaTeX]TeX}, numbers=none, frame=none,
                   backgroundcolor=\color{white}, xleftmargin=0pt},
  fontlower=\small\linespread{1}\selectfont,
}

% Apenas o código LaTeX (usado para figuras, tabelas e outros objetos
% flutuantes, cujo resultado aparece no próprio texto do exemplo).
\newtcblisting{exemplocodigo}{
  enhanced, listing only,
  colback=white, colframe=exemploborda, boxrule=0.4pt, arc=1pt,
  left=4pt, right=4pt, top=2pt, bottom=2pt,
  title={Código LaTeX}, coltitle=black, colbacktitle=black!8,
  fonttitle=\sffamily\footnotesize,
  listing options={language={[LaTeX]TeX}, numbers=none, frame=none,
                   backgroundcolor=\color{white}, xleftmargin=0pt},
}
